\chapter{SQL et données}\label{ch:iv3:sql-and-data}

\begin{ivtrains}
Écrire une requête correcte sur les lignes délicates, et lire une requête qui semble juste sans l'être~: jointures qui perdent ou multiplient des lignes, les trois valeurs de la logique SQL, fenêtres de calcul, jointures as-of, trous et mauvaises données. Chaque table est imprimée, de sorte que l'on peut exécuter la requête de tête. Un énoncé précise le dialecte lorsque la réponse en dépend.
\end{ivtrains}

\begin{question}[$\star$ \iqroles{QD, Dev, DS, QR} \iqfirm{toutes} \iqlevel{junior}]\label{iq:iv3:sql-and-data:1}
Voici une table \texttt{trades}, où \texttt{price} est un double et \texttt{qty} un entier. Écrivez une requête donnant, pour chaque symbole et chaque jour, le prix moyen pondéré par les volumes, c'est-à-dire la somme des prix multipliés par les quantités, divisée par la somme des quantités. Ne retenez que les couples symbole-jour comptant plus de 2 transactions. Que renvoie-t-elle ici~?
\begin{center}\small
\begin{tabular}{llrr}
symbol & day & price & qty\\\hline
AAA & 2026-03-02 & 100.0 & 200\\
AAA & 2026-03-02 & 102.0 & 100\\
AAA & 2026-03-02 & 101.0 & 100\\
BBB & 2026-03-02 & 50.0 & 400\\
BBB & 2026-03-02 & 52.0 & 100\\
AAA & 2026-03-03 & 104.0 & 100
\end{tabular}
\end{center}
\end{question}

\begin{question}[$\star$ \iqroles{QD, Dev, DS, QR} \iqfirm{toutes} \iqlevel{junior}]\label{iq:iv3:sql-and-data:2}
La table de gauche est \texttt{orders} et celle de droite est \texttt{fills}. Vous voulez chaque ordre, avec la quantité exécutée sur la plateforme X, ou NULL s'il n'en a aucune. Que renvoie chaque requête du \cref{lst:iv3:sq:venue}, et laquelle est correcte~?
\begin{center}\small
\begin{tabular}{rl}
order\_id & symbol\\\hline
1 & AAA\\
2 & BBB\\
3 & CCC
\end{tabular}\qquad
\begin{tabular}{rlr}
order\_id & venue & qty\\\hline
1 & X & 100\\
1 & Y & 50\\
2 & Y & 70
\end{tabular}
\end{center}
\omcode{code/interviews-3/07-sql-and-data/sql/venue_fills.sql}{1}{8}{Deux jointures à gauche.}\label{lst:iv3:sq:venue}
\end{question}

\begin{question}[$\star$ \iqroles{QD, Dev, DS} \iqfirm{toutes} \iqlevel{junior}]\label{iq:iv3:sql-and-data:3}
La table de gauche est \texttt{orders} et celle de droite est \texttt{blocked}. Que renvoie le \cref{lst:iv3:sq:notin}, et pourquoi~? Réécrivez-la de sorte qu'elle renvoie chaque ordre dont le trader ne figure pas dans la liste des traders bloqués. L'ordre 3 doit-il figurer dans votre réponse~?
\begin{center}\small
\begin{tabular}{rl}
order\_id & trader\\\hline
1 & ann\\
2 & bob\\
3 & NULL
\end{tabular}\qquad
\begin{tabular}{l}
trader\\\hline
bob\\
NULL
\end{tabular}
\end{center}
\omcode{code/interviews-3/07-sql-and-data/sql/not_in.sql}{1}{3}{Ordres émis par des traders non bloqués.}\label{lst:iv3:sq:notin}
\end{question}

\begin{question}[$\star$ \iqroles{QD, Dev, DS} \iqfirm{toutes} \iqlevel{junior}]\label{iq:iv3:sql-and-data:4}
Une table \texttt{quotes} possède une colonne \texttt{bid} de type double, qui contient les cinq valeurs 10.0, NULL, 20.0, 20.0, NULL. Que renvoient, valeur par valeur, les deux instructions du \cref{lst:iv3:sq:counts}~? Laquelle des deux dernières colonnes de la première instruction est le cours acheteur moyen~?
\omcode{code/interviews-3/07-sql-and-data/sql/counts.sql}{1}{4}{Compter en présence de NULL.}\label{lst:iv3:sq:counts}
\end{question}

\begin{question}[$\star$ \iqroles{QD, Dev, DS} \iqfirm{toutes} \iqlevel{junior}]\label{iq:iv3:sql-and-data:5}
Une table \texttt{fills} compte deux lignes, avec les colonnes entières \texttt{filled} et \texttt{ordered} contenant (30, 100) et (45, 50). Que renvoie le \cref{lst:iv3:sq:ratio} sous SQLite, sous DuckDB et sous PostgreSQL~? Écrivez une version qui renvoie le même résultat sur les trois.
\omcode{code/interviews-3/07-sql-and-data/sql/fill_ratio.sql}{1}{1}{La part de l'ordre exécutée.}\label{lst:iv3:sq:ratio}
\end{question}

\begin{question}[$\star$ \iqroles{QD, Dev, DS} \iqfirm{toutes} \iqlevel{junior}]\label{iq:iv3:sql-and-data:6}
Un flux renvoie parfois une transaction une seconde fois, et la corrige parfois en envoyant le même \texttt{trade\_id} avec une nouvelle quantité. Dans la table \texttt{trades} ci-dessous, \texttt{seq} est l'ordre d'arrivée. Listez les identifiants qui se répètent, et dites si leurs exemplaires concordent. Donnez ensuite la quantité totale par symbole, en retenant la dernière version de chaque transaction. \texttt{SELECT DISTINCT} aurait-il donné les mêmes totaux~?
\begin{center}\small
\begin{tabular}{rrlr}
seq & trade\_id & symbol & qty\\\hline
1 & 1 & AAA & 100\\
2 & 2 & AAA & 200\\
3 & 2 & AAA & 200\\
4 & 3 & BBB & 50\\
5 & 3 & BBB & 60\\
6 & 4 & BBB & 70
\end{tabular}
\end{center}
\end{question}

\begin{question}[$\star\star$ \iqroles{QD, Dev, DS, QR} \iqfirm{toutes} \iqlevel{junior, confirmé}]\label{iq:iv3:sql-and-data:7}
Voici une table \texttt{prices}. Que fait le \cref{lst:iv3:sq:bare} sous SQLite, sous DuckDB et sous PostgreSQL~? Écrivez une requête donnant le dernier prix de chaque symbole, correcte sur les trois.
\begin{center}\small
\begin{tabular}{lrr}
symbol & ts & price\\\hline
AAA & 1 & 100.0\\
AAA & 3 & 101.0\\
AAA & 2 & 99.0\\
BBB & 5 & 50.0\\
BBB & 4 & 51.0
\end{tabular}
\end{center}
\omcode{code/interviews-3/07-sql-and-data/sql/latest_bare.sql}{1}{3}{Le dernier prix, peut-être.}\label{lst:iv3:sq:bare}
\end{question}

\begin{question}[$\star\star$ \iqroles{QD, Dev, DS} \iqfirm{toutes} \iqlevel{junior, confirmé}]\label{iq:iv3:sql-and-data:8}
La table \texttt{pnl} compte une ligne par trader. On vous demande les 2 meilleurs traders par profit sur chaque desk. Exécutez de tête \texttt{ROW\_NUMBER}, \texttt{RANK} et \texttt{DENSE\_RANK} sur \texttt{(PARTITION BY desk ORDER BY pnl DESC)}~: que donne chacune sur cette table~? Quels traders un filtre à 2 ou moins retient-il avec chacune d'elles~? Laquelle utiliseriez-vous, et que devez-vous demander en premier~?
\begin{center}\small
\begin{tabular}{llr}
desk & trader & pnl\\\hline
rates & ann & 50\\
rates & bob & 50\\
rates & cid & 40\\
rates & dan & 30\\
fx & eve & 70\\
fx & fay & 60\\
fx & gus & 60
\end{tabular}
\end{center}
\end{question}

\begin{question}[$\star\star$ \iqroles{QD, Dev, DS, QT} \iqfirm{toutes} \iqlevel{junior, confirmé}]\label{iq:iv3:sql-and-data:9}
La table \texttt{fills} contient les quantités signées d'un seul symbole, où \texttt{seq} est l'ordre dans lequel les exécutions sont survenues, et deux exécutions peuvent partager un même horodatage \texttt{ts}. Une position courte est contraire au règlement, donc la requête du \cref{lst:iv3:sq:position} vérifie l'absence de position cumulée négative. Qu'affiche-t-elle, et détecte-t-elle l'infraction~? Corrigez-la.
\begin{center}\small
\begin{tabular}{rrr}
seq & ts & qty\\\hline
1 & 10 & 100\\
2 & 20 & $-150$\\
3 & 20 & 100\\
4 & 30 & $-50$
\end{tabular}
\end{center}
\omcode{code/interviews-3/07-sql-and-data/sql/position.sql}{1}{3}{Une position cumulée.}\label{lst:iv3:sq:position}
\end{question}

\begin{question}[$\star\star$ \iqroles{QD, Dev, DS} \iqfirm{toutes} \iqlevel{junior, confirmé}]\label{iq:iv3:sql-and-data:10}
La table de gauche est \texttt{orders} et celle de droite est \texttt{fills}. Le \Cref{lst:iv3:sq:fanout} doit donner, pour chaque desk, la quantité totale commandée et le total exécuté. Qu'affiche-t-il, qu'est-ce qui ne va pas dans chacun de ces deux nombres, et comment le corrigez-vous~?
\begin{center}\small
\begin{tabular}{rlr}
order\_id & desk & qty\\\hline
1 & rates & 100\\
2 & rates & 50\\
3 & fx & 200
\end{tabular}\qquad
\begin{tabular}{rr}
order\_id & filled\\\hline
1 & 60\\
1 & 40\\
3 & 150
\end{tabular}
\end{center}
\omcode{code/interviews-3/07-sql-and-data/sql/fanout.sql}{1}{4}{Commandé et exécuté, par desk.}\label{lst:iv3:sq:fanout}
\end{question}

\begin{question}[$\star\star$ \iqroles{QD, Dev, DS, QR} \iqfirm{toutes} \iqlevel{junior, confirmé}]\label{iq:iv3:sql-and-data:11}
Dans la table \texttt{closes}, \texttt{day} est un numéro de jour calendaire, et les jours 4 et 5 forment un week-end sans ligne. Quelles sont, ligne par ligne, les deux colonnes de moyennes produites par le \cref{lst:iv3:sq:moving}~? Laquelle est une moyenne mobile sur trois jours~?
\begin{center}\small
\begin{tabular}{rr}
day & price\\\hline
1 & 10.0\\
2 & 20.0\\
3 & 30.0\\
6 & 60.0\\
7 & 90.0
\end{tabular}
\end{center}
\omcode{code/interviews-3/07-sql-and-data/sql/moving.sql}{1}{5}{Deux moyennes mobiles.}\label{lst:iv3:sq:moving}
\end{question}

\begin{question}[$\star\star$ \iqroles{QD, Dev, DS, QR} \iqfirm{toutes} \iqlevel{junior, confirmé}]\label{iq:iv3:sql-and-data:12}
Le \Cref{lst:iv3:sq:barsbad} devrait transformer la table \texttt{ticks} en une barre par minute, avec l'ouverture, le plus haut, le plus bas et la clôture. Combien de lignes renvoie-t-il, et quelles sont-elles~? Pourquoi~? Corrigez-le.
\begin{center}\small
\begin{tabular}{rrr}
minute & seq & price\\\hline
1 & 1 & 100.0\\
1 & 2 & 103.0\\
1 & 3 & 101.0\\
2 & 4 & 102.0\\
2 & 5 & 99.0
\end{tabular}
\end{center}
\omcode{code/interviews-3/07-sql-and-data/sql/bars_bad.sql}{1}{8}{Barres à partir des ticks.}\label{lst:iv3:sq:barsbad}
\end{question}

\begin{question}[$\star\star$ \iqroles{QD, Dev, DS, QR} \iqfirm{toutes} \iqlevel{junior, confirmé}]\label{iq:iv3:sql-and-data:13}
La table de gauche est \texttt{trades} et celle de droite est \texttt{quotes}, avec \texttt{ts} en secondes. Pour chaque transaction, renvoyez le cours acheteur et le cours vendeur de la dernière cotation survenue à l'instant de la transaction ou avant, mais traitez comme manquante une cotation vieille de plus de 10 secondes. Utilisez une sous-requête corrélée, sans aucune syntaxe de jointure particulière. Que renvoyez-vous pour chaque transaction ici~? Quel est le coût de votre requête avec et sans index sur \texttt{quotes.ts}~?
\begin{center}\small
\begin{tabular}{rr}
trade\_id & ts\\\hline
1 & 2\\
2 & 5\\
3 & 12\\
4 & 30
\end{tabular}\qquad
\begin{tabular}{rrr}
ts & bid & ask\\\hline
3 & 99.9 & 100.1\\
10 & 100.3 & 100.6\\
12 & 100.4 & 100.7
\end{tabular}
\end{center}
\end{question}

\begin{question}[$\star\star$ \iqroles{QD, Dev, LLE, DS} \iqfirm{teneur de marché, crypto, hedge fund} \iqlevel{junior, confirmé}]\label{iq:iv3:sql-and-data:14}
Une table \texttt{msgs} contient une seule colonne entière, \texttt{seq}, qui porte les numéros de séquence reçus d'un flux censé ne comporter aucun trou. Ses lignes sont 1, 2, 3, 6, 7 et 11. Écrivez une requête qui renvoie chaque plage manquante sous la forme d'un premier et d'un dernier numéro. Une requête plus courte, le plus grand moins le plus petit, plus un, moins le nombre de lignes, donne le nombre de numéros manquants~: pourquoi ne suffit-elle pas~? Que ne peut jamais vous dire aucune requête sur cette table~?
\end{question}

\begin{question}[$\star\star$ \iqroles{QD, Dev, DS, QR} \iqfirm{teneur de marché, crypto, hedge fund} \iqlevel{junior, confirmé}]\label{iq:iv3:sql-and-data:15}
La table \texttt{quotes} ci-dessous contient un seul symbole, dans l'ordre d'arrivée \texttt{seq}. Écrivez une requête qui renvoie chaque ligne suspecte avec une seule raison, en vérifiant dans cet ordre~: un cours acheteur supérieur ou égal au cours vendeur~; une taille nulle ou négative~; un horodatage qui ne croît pas~; et un prix milieu ayant varié de plus de 5~\% par rapport à celui de la ligne précédente. Quelles lignes sont renvoyées~? Laquelle d'entre elles est une bonne ligne~? Comment empêcheriez-vous qu'une mauvaise cotation ne signale deux lignes~?
\begin{center}\small
\begin{tabular}{rrrrr}
seq & ts & bid & ask & bid\_size\\\hline
1 & 1 & 100.0 & 100.2 & 500\\
2 & 2 & 100.1 & 100.3 & 400\\
3 & 3 & 100.4 & 100.2 & 300\\
4 & 4 & 100.1 & 100.3 & 0\\
5 & 4 & 100.1 & 100.3 & 200\\
6 & 6 & 110.0 & 110.2 & 100\\
7 & 7 & 100.2 & 100.4 & 300
\end{tabular}
\end{center}
\end{question}

\begin{question}[$\star\star$ \iqroles{QD, Dev, LLE, DS} \iqfirm{toutes} \iqlevel{junior, confirmé}]\label{iq:iv3:sql-and-data:16}
Une table \texttt{latency} contient une seule colonne entière \texttt{us}, qui porte les valeurs 10, 12, 13, 15, 20, 22, 25, 30, 80 et 200. Sur DuckDB, que renvoient \texttt{quantile\_disc(us, 0.9)} et \texttt{quantile\_cont(us, 0.9)}, et pourquoi diffèrent-elles~? SQLite 3.37 n'en possède aucune des deux, donc calculez-y le 90\up{e}~centile selon la règle du rang le plus proche~: la plus petite valeur telle qu'au moins 90~\% des lignes lui soient inférieures ou égales. Quelle définition retiendriez-vous pour la latence~?
\end{question}

\begin{question}[$\star\star$ \iqroles{QD, Dev, DS} \iqfirm{toutes} \iqlevel{junior, confirmé}]\label{iq:iv3:sql-and-data:17}
Sur SQLite, une table \texttt{trades} possède l'index du \cref{lst:iv3:sq:plans}, sur \texttt{(symbol, ts)}. Pour chacune des cinq requêtes, le plan recherche-t-il dans l'index ou parcourt-il la table entière, et a-t-il besoin d'un tri séparé~? Expliquez la règle qui sous-tend vos réponses.
\omcode{code/interviews-3/07-sql-and-data/sql/index_plans.sql}{1}{13}{Un index, cinq requêtes.}\label{lst:iv3:sq:plans}
\end{question}

\begin{question}[$\star\star\star$ \iqroles{QD, Dev, DS, QR} \iqfirm{toutes} \iqlevel{junior, confirmé}]\label{iq:iv3:sql-and-data:18}
La table \texttt{daily} contient une ligne par jour de négociation, où \texttt{day} croît sans être forcément consécutif. Écrivez une requête qui renvoie le premier jour, le dernier jour et la longueur de la plus longue suite de lignes consécutives dont le profit est strictement positif. En cas d'égalité, retenez la plus ancienne. Que renvoie-t-elle ici~? N'utilisez ni boucle ni récursion.
\begin{center}\small
\begin{tabular}{rr}
jour & pnl\\\hline
1 & 5\\
2 & 3\\
3 & $-2$\\
6 & 1\\
7 & 2\\
8 & 4\\
9 & $-1$\\
10 & 7
\end{tabular}
\end{center}
\end{question}

\begin{question}[$\star\star\star$ \iqroles{QD, Dev, DS} \iqfirm{toutes} \iqlevel{junior, confirmé}]\label{iq:iv3:sql-and-data:19}
La table \texttt{events} contient les actions d'un trader, avec \texttt{ts} exprimé en secondes. Une session est une suite d'événements d'un même trader dans laquelle chaque événement suit le précédent au plus 300 secondes plus tard, un écart plus long ouvrant une nouvelle session. Écrivez une requête qui renvoie, pour chaque session, le trader, le numéro de session, le premier et le dernier instant et le nombre d'événements. Que renvoie-t-elle ici~? Pourquoi un écart de 300 secondes exactement est-il le cas à tester~?
\begin{center}\small
\begin{tabular}{lr}
trader & ts\\\hline
ann & 100\\
ann & 200\\
ann & 700\\
ann & 800\\
ann & 1500\\
bob & 150\\
bob & 460\\
bob & 760
\end{tabular}
\end{center}
\end{question}

\begin{question}[$\star\star\star$ \iqroles{QD, Dev, DS, QR} \iqfirm{toutes} \iqlevel{junior, confirmé}]\label{iq:iv3:sql-and-data:20}
Dans la table \texttt{closes}, un prix NULL signifie qu'il n'y a pas eu de transaction ce jour-là. Remplissez chaque NULL avec le dernier prix connu, et laissez tel quel un NULL qui n'a rien avant lui. Écrivez-la pour SQLite 3.37, qui ne prend pas en charge \texttt{IGNORE NULLS}, et expliquez comment DuckDB peut le faire en une seule ligne. Que renvoyez-vous ici~? Pourquoi ne rempliriez-vous pas aussi vers l'arrière~?
\begin{center}\small
\begin{tabular}{rr}
jour & prix\\\hline
1 & NULL\\
2 & 10.0\\
3 & NULL\\
4 & NULL\\
5 & 12.0\\
6 & NULL
\end{tabular}
\end{center}
\end{question}

\begin{question}[$\star\star\star$ \iqroles{QD, Dev, DS, QR} \iqfirm{teneur de marché, hedge fund, crypto} \iqlevel{confirmé}]\label{iq:iv3:sql-and-data:21}
La table de gauche est \texttt{trades} et celle de droite \texttt{quotes}, chacune comptant des millions de lignes par symbole, et aucune des deux ne possède d'index. Pour chaque transaction, renvoyez le \texttt{mid} de la dernière cotation du même symbole survenue à l'instant de la transaction ou avant, ou NULL s'il n'y en a aucune. Une cotation au même instant qu'une transaction compte. Donnez la requête sur DuckDB, en utilisant sa jointure spéciale. Donnez ensuite une requête sans cette jointure ni sous-requête corrélée, qui s'exécute sur SQLite 3.37. Quel est le coût de chacune, et que renvoyez-vous ici~?
\begin{center}\small
\begin{tabular}{lr}
symbole & ts\\\hline
AAA & 5\\
AAA & 12\\
BBB & 1\\
BBB & 7
\end{tabular}\qquad
\begin{tabular}{lrr}
symbole & ts & mid\\\hline
AAA & 3 & 100.0\\
AAA & 12 & 101.0\\
BBB & 2 & 50.0\\
BBB & 9 & 51.0
\end{tabular}
\end{center}
\end{question}

\begin{question}[$\star\star\star$ \iqroles{QD, Dev, DS, QR} \iqfirm{hedge fund, banque} \iqlevel{confirmé}]\label{iq:iv3:sql-and-data:22}
Un ticker peut passer d'une société à une autre. La table de gauche, \texttt{master}, donne le ticker de chaque instrument et les dates pendant lesquelles l'instrument le portait, de \texttt{valid\_from} jusqu'à \texttt{valid\_to} non compris, et la table de droite est \texttt{trades}. Les dates sont en texte ISO. Le \Cref{lst:iv3:sq:pit} est une jointure que quelqu'un a écrite. Que renvoie-t-elle, et quels sont les deux problèmes distincts de ce résultat~? Corrigez la requête, et écrivez une requête de qualité des données qui repère les mauvaises lignes de \texttt{master}.
\begin{center}\small
\begin{tabular}{rlll}
id\_instrument & ticker & valide\_du & valide\_au\\\hline
1 & ABC & 2020-01-01 & 2024-05-01\\
2 & ABC & 2024-05-01 & 9999-12-31\\
3 & XYZ & 2020-01-01 & 2023-01-01\\
4 & XYZ & 2022-06-01 & 9999-12-31
\end{tabular}\qquad
\begin{tabular}{llr}
ticker & jour & qty\\\hline
ABC & 2024-04-30 & 10\\
ABC & 2024-05-01 & 20\\
XYZ & 2022-07-01 & 30
\end{tabular}
\end{center}
\omcode{code/interviews-3/07-sql-and-data/sql/pit_between.sql}{1}{4}{Une jointure sur les dates de validité.}\label{lst:iv3:sq:pit}
\end{question}

\begin{question}[$\star\star\star$ \iqroles{QD, Dev, DS, QR} \iqfirm{teneur de marché, hedge fund, crypto} \iqlevel{confirmé}]\label{iq:iv3:sql-and-data:23}
La table \texttt{trades} contient un temps entier \texttt{ts} en secondes, un réel double précision \texttt{price} et un entier \texttt{qty}. Pour chaque transaction, donnez le prix moyen pondéré par les volumes (VWAP) sur les cinq minutes se terminant à cette transaction, c'est-à-dire sur les transactions telles que $t - 300 < \text{ts} \le t$. Utilisez une fonction de fenêtre avec un cadre en temps, et non en lignes. Que renvoyez-vous ici~? Deux transactions partagent l'instant 420~: se voient-elles mutuellement~? Que faudrait-il changer si \texttt{ts} était un réel~?
\begin{center}\small
\begin{tabular}{rrr}
ts & prix & qty\\\hline
0 & 100.0 & 10\\
120 & 101.0 & 10\\
299 & 102.0 & 20\\
300 & 103.0 & 10\\
420 & 100.0 & 10\\
420 & 99.0 & 30
\end{tabular}
\end{center}
\end{question}

\begin{question}[$\star\star\star\star$ \iqroles{QD, Dev, DS, QT} \iqfirm{toutes} \iqlevel{confirmé}]\label{iq:iv3:sql-and-data:24}
La table \texttt{fills} contient les achats et les ventes d'un même symbole, dans l'ordre \texttt{seq}, et la position ne devient jamais négative. Les ventes sont appariées aux achats selon la règle premier entré, premier sorti~: chaque vente épuise les actions les plus anciennes encore détenues. Dans une seule instruction SQL, avec des fonctions de fenêtre mais sans boucle ni récursion, renvoyez une ligne par paire achat-vente appariée, avec la quantité et le profit, sur SQLite 3.37 et sur DuckDB. Quel est le profit total ici~? Votre requête tient-elle encore si la position peut devenir vendeuse~?
\begin{center}\small
\begin{tabular}{rlrr}
seq & sens & prix & qty\\\hline
1 & B & 100 & 100\\
2 & B & 110 & 50\\
3 & S & 120 & 120\\
4 & B & 90 & 30\\
5 & S & 130 & 60
\end{tabular}
\end{center}
\end{question}

\begin{question}[$\star\star\star\star$ \iqroles{QD, Dev, DS} \iqfirm{teneur de marché, crypto} \iqlevel{confirmé}]\label{iq:iv3:sql-and-data:25}
La table \texttt{events} est l'historique du buy side d'un même carnet d'ordres~: à l'étape \texttt{seq}, la quantité au repos à \texttt{price} varie de \texttt{delta}. Dans une seule instruction SQL, sur SQLite 3.37 et sur DuckDB, renvoyez pour chaque étape le meilleur cours acheteur après cette étape, c'est-à-dire le prix le plus élevé ayant une quantité positive, ainsi que la quantité à ce prix, ou NULL si ce côté est vide. Que renvoyez-vous ici~? Quel est le coût de votre requête, et la lanceriez-vous ainsi sur une journée de données~?
\begin{center}\small
\begin{tabular}{rrr}
seq & prix & delta\\\hline
1 & 100 & 50\\
2 & 101 & 20\\
3 & 100 & 30\\
4 & 101 & $-20$\\
5 & 99 & 40\\
6 & 100 & $-80$
\end{tabular}
\end{center}
\end{question}
